\section{The Solar Sign, the Early Discrepancy, and the Church's Response}

\subsection{What the contemporary record supports}

Avelino de Almeida's report in \emph{O Século}, published 15 October, is important because it was near-event, public, and not written as an ecclesial brief. He described a rain-soaked journey and crowd; observers he called dispassionate estimated thirty to forty thousand. Later Shrine presentations often give fifty to seventy thousand. No controlled count exists, so a range is more honest than a talismanic number.

After the rain stopped and the cloud cover broke, Almeida saw the solar disc as a dull silver plate that could be viewed without the ordinary glare. A great cry rose. He described trembling or dancing, while recording that participants offered different reports: rotation, color changes, an apparent descent, and, for some, religious images. He did not say that every person had an identical perception. Other contemporary and later witnesses supplied accounts from the Cova and surrounding distances, with both convergence and variation.

At least five propositions can therefore be separated:

\begin{enumerate}
  \item A large crowd gathered at a time and place announced in advance.
  \item The day was wet, and a visible change in sky conditions occurred near the expected period.
  \item Many witnesses, including persons not disposed to credulity, reported a striking solar or atmospheric appearance.
  \item Descriptions were not uniform, and the event was not observed as a global displacement of the astronomical sun.
  \item Historical testimony alone does not settle whether the local event was miraculous in cause, providentially timed natural phenomena, physiological effects, collective expectation acting upon ambiguous stimuli, or some combination.
\end{enumerate}

The 1930 pastoral letter itself noted that no astronomical observatory registered a solar movement. The bishop used that absence to argue that the phenomenon was not an ordinary movement of the sun and received it within the whole case as supernatural. A contemporary physical inference is narrower: the star did not leave its astronomical course. Had it done so, the consequences and observation would not have been local. This does not explain away what the crowd saw; it prevents the devotional phrase “the sun fell” from being mistaken for astrophysics.

Natural proposals---thin cloud, atmospheric optics, retinal afterimages from staring at the sun, suggestion, and variable local conditions---can explain some features. They do not automatically explain the announced timing, the scale of response, or every testimony. “Miracle” likewise should not become a placeholder for ignorance. Catholic judgment considers the sign within persons, message, prediction, fruits, and doctrine; it need not furnish a meteorological mechanism.

\interpretivekey{sign rather than spectacle}{In Scripture a sign directs beyond itself. The crowd's experience serves Fátima only if it redirects attention from the sky to repentance, Christ, and the life of the Church. Repeated attempts to stare at the sun are dangerous, unnecessary, and spiritually perverse. Faith does not require reenactment.}

\subsection{The statement that the war would end “today”}

The same near-event record that strengthens the historical case also preserves its hardest verbal discrepancy. On 16 October the parish priest recorded Lúcia as reporting: \latin{a guerra acaba ainda hoje}---“the war ends today”---and that the soldiers would return soon. Formigão questioned her on 19 October because newspapers still reported combat. She answered that she did not know; she had heard it end on the thirteenth. When reminded that an earlier answer may have meant “soon,” she allowed that she might not remember or might not have understood the Lady correctly. Other early testimony gives “soon” or a conditional relation to amendment. The later Fourth Memoir says that the war “is going to end.” The Armistice came on 11 November 1918.

Four inadequate responses should be refused.

\begin{itemize}
  \item \textbf{Silent harmonization} replaces the 1917 document with the later wording and deprives the reader of evidence.
  \item \textbf{Ad hoc fulfillment} invents an invisible sense in which the war ended that day, without historical warrant.
  \item \textbf{Total collapse} treats one difficulty as though it erased every distinct witness, fruit, and ecclesial judgment.
  \item \textbf{Verbal inerrancy} assumes that approval of an interior vision guarantees perfect reception, memory, and expression of every sentence.
\end{itemize}

Cardinal Ratzinger's 2000 anthropology of vision provides a principled account rather than an excuse invented for this phrase. In interior vision, the recipient is touched by a reality beyond ordinary sense, but sees within the modes of representation and capacities available to the person. The image or word is received through a human subject, not downloaded around that subject. Even ordinary exterior perception includes translation; interior perception involves it more deeply. The child may hear, understand, remember, or report imperfectly without the entire experience becoming fraudulent.

The bishop knew the difficulty. The critical documentation's introduction notes that the delayed end of the war challenged credibility and that the commission examined proposed resolutions. His 1930 act nevertheless judged the visions as a whole worthy of belief. That judgment deserves respect in its actual form: neither an erasure of the discrepancy nor a demand that Catholics solve it.

\subsection{From spontaneous pilgrimage to formal inquiry}

The response developed in overlapping lines rather than by one institutional decision.

\begin{fourstable}{Period}{Popular and civil response}{Ecclesial response}{Meaning}
1917 & Crowds grow; press attention spreads; children are questioned and detained; money and petitions accumulate. & Parish clergy record testimony while maintaining caution; the diocese is still being restored. & The event becomes public before the hierarchy judges it.\\
1918--1921 & Pilgrimage continues; a small chapel is built in 1919; the first Mass there is celebrated in 1921. & Bishop Correia da Silva assumes the restored see and begins governing the movement. & Devotion has facts and risks requiring pastoral order.\\
1922--1924 & The chapel is dynamited in March 1922, rebuilt, and reopened in January 1923. & The bishop appoints a commission in 1922; Lúcia gives a formal examination in 1924. & Opposition and enthusiasm alike remain subordinate to evidence.\\
1925--1930 & Pilgrimage and cult consolidate within Portugal and beyond. & Commission report is reviewed; bishop issues the 13 Oct. 1930 judgment and cult permission. & A thirteen-year public history is bounded by a competent act.\\
\end{fourstable}

The chapel was constructed between 28 April and 15 June 1919. Its bombing on 6 March 1922 partially destroyed it; restoration followed. Persecution is not automatic proof of divine origin---false claims can be opposed---but the episode shows that the shrine did not emerge from uncontested clerical manufacture. Nor was popular persistence by itself decisive. The bishop's letter reasons from the children's character, doctrinal content, associated phenomena, persecution, fruits, and the commission's work, then submits his judgment to the Holy See.

The act's two clauses are deliberately modest. The first addresses credibility of the visions. The second permits cult. The message thereafter gained a stable ecclesial home in which sacraments, preaching, care of pilgrims, archive, and scrutiny could shape popular devotion. Institutionalization can obscure origins; it can also protect a charism from private ownership.

\subsection{A sign judged inside a whole}

The solar event cannot carry Fátima alone. Without message and holiness it would be a curiosity. Without contemporary testimony the narrative would be historically thinner. Without ecclesial judgment it would remain a disputed public claim. Without conversion its religious fruit would fail.

The Catholic conclusion is proportionate. One may rationally take the mass testimony and announced timing as powerful support within the approved whole. One may remain uncertain about mechanism. One may not assert that astrophysics was suspended globally, that every observer saw the same thing, or that doubt about an optical detail is rejection of the faith. The faith rests on Christ crucified and risen; the Fátima sign, if received, points there.
